\uuid{Lq5K}
\titre{Changement de variables}
\auteur{hanani ; mzari}
\organisation{exo7}
\licence{CC BY-NC-SA 4.0}
\datecreate{2020-01-20}
\niveau{L1}
\module{Analyse}
\chapitre{Calcul d'intégrales}
\sousChapitre{Primitives diverses}
\difficulte{3}
\theme{Primitives, changement de variables}
\exo7qcmid{lille-2.primitives.032}

\contenu{
\begin{question}
\qclassification{127.11}
\qtags{motcle=changement de variables}
On se place sur $]0,+\infty[$. Le changement de variable $u=\mathrm{e}^{-x}$ donne
\begin{answers}
\bad{$\displaystyle \int \frac{\mathrm{d}x}{\sqrt{\mathrm{e}^{2x}-1}}=\int \frac{\mathrm{d}u}{\sqrt{1-u^2}}$.}
\good{$\displaystyle \int \frac{\mathrm{d}x}{\sqrt{\mathrm{e}^{2x}-1}}=-\arcsin (u)+k$, $k\in \Rr$.}
\bad{$\displaystyle \int \frac{\mathrm{e}^{-x}\mathrm{d}x}{\sqrt{\mathrm{e}^{2x}-1}}=\int \frac{u\mathrm{d}u}{\sqrt{1-u^2}}$.}
\good{$\displaystyle \int \frac{\mathrm{e}^{-x}\mathrm{d}x}{\sqrt{\mathrm{e}^{2x}-1}}=\sqrt{(1-u^2)}+k$, $k\in \Rr$.}
\end{answers}
\vskip2mm
\begin{explanations}
Avec $u=\mathrm{e}^{-x}$, on aura : $\mathrm{d}u=-\mathrm{e}^{-x}\mathrm{d}x$. D'où
$$\displaystyle \int \frac{\mathrm{d}x}{\sqrt{\mathrm{e}^{2x}-1}}=\int \frac{\mathrm{e}^{-x}\mathrm{d}x}{\sqrt{1-\mathrm{e}^{-2x}}}=\int \frac{-\mathrm{d}u}{\sqrt{1-u^2}}=-\arcsin(u)+k,\; k\in \Rr.$$
Et
$$\displaystyle \int \frac{\mathrm{e}^{-x}\mathrm{d}x}{\sqrt{\mathrm{e}^{2x}-1}}=\int \frac{\mathrm{e}^{-2x}\mathrm{d}x}{\sqrt{1-\mathrm{e}^{-2x}}}=\int \frac{-u\mathrm{d}u}{\sqrt{1-u^2}}=\sqrt{(1-u^2)}+k,\; k\in \Rr.$$
\end{explanations}
\end{question}
}
